% Section content template for individual Educative lessons
% This file contains the actual content for a specific section
% Generated from Educative API section content

% Section: Quiz on YouTube's Design
% Section ID: 6669179066777600
% Section Slug: quiz-on-youtubes-design
% Generated: 2025-12-07T22:38:51.538232

The following quiz includes questions about the Netflix system to evaluate your understanding of a streaming platform. Since YouTube and Netflix are streaming platforms, they share similar challenges when designing these systems. By exploring Netflix's design, we can gain valuable insights that can be applied to YouTube's architecture as well.
\\
Let's dive in!

\subsection*{Quiz}
\vspace{0.3cm}
\noindent
\textbf{Question 1:} Which pair of requirements cannot be compromised when designing Netflix?
\\[0.2cm]
\begin{itemize}
 \item \textbf{[CORRECT]} Low latency and high availability
\\[0.1cm]
\textit{Explanation:} 
Streaming services highly require these two design goals to provide a good QoE to end-user.
 \item Low latency and HD video
\\[0.1cm]
\textit{Explanation:} 
It is extremely important for the service to be available even if the video is not high quality
 \item High availability and HD video
\\[0.1cm]
\textit{Explanation:} 
This could lead to buffering and lag issues. Therefore, low latency is required.
 \item Consistency and HD video
\\[0.1cm]
\textit{Explanation:} 
Consistency isn't the most challenging or desired metric when it comes to streaming services.
\end{itemize}
\vspace{0.3cm}
\noindent
\textbf{Question 2:} Consider the following design to upload content to the Netflix service:
\\
\pandocbounded{\includegraphics[keepaspectratio]{/api/collection/10370001/4941429335392256/page/6669179066777600/image/5945405562290176?page_type=collection_lesson&get_optimised=true}}
\\
Note that there are two unnamed components. Select the correct option to complete the design.
\\[0.2cm]
\begin{itemize}
 \item \includegraphics[width=\linewidth,height=1.5625in,keepaspectratio]{/api/collection/10370001/4941429335392256/page/6669179066777600/image/4654959679635456?page_type=collection_lesson&get_optimised=true
}
\\[0.1cm]
\textit{Explanation:} 
Metadata storage cannot be the first point of contact for the end-users.
 \item \includegraphics[width=\linewidth,height=1.5625in,keepaspectratio]{/api/collection/10370001/4941429335392256/page/6669179066777600/image/6430911476662272?page_type=collection_lesson&get_optimised=true
}
\\[0.1cm]
\textit{Explanation:} 
It doesn't make sense for the uploading servers and Mapreduce to store data in colocation sites instead of a datastore.
 \item \textbf{[CORRECT]} \includegraphics[width=\linewidth,height=1.5625in,keepaspectratio]{/api/collection/10370001/4941429335392256/page/6669179066777600/image/5282981189386240?page_type=collection_lesson&get_optimised=true
}
\\[0.1cm]
\textit{Explanation:} 
Users' requests are distributed among uploading servers using load balancers. Also , we can use a datastore to save (meta)data from uploading servers and Mapreduce.
 \item \includegraphics[width=\linewidth,height=1.5625in,keepaspectratio]{/api/collection/10370001/4941429335392256/page/6669179066777600/image/6142871709548544?page_type=collection_lesson&get_optimised=true}
\\[0.1cm]
\textit{Explanation:} 
Web servers require load balancers before them in practical designs. Moreover , uploading servers would extract metadata from raw videos and store it in a database that isn't a key-value store.
\end{itemize}
\vspace{0.3cm}
\noindent
\textbf{Question 3:} Which encoding technique is the most effective in terms of streaming to end users?
\\[0.2cm]
\begin{itemize}
 \item Use a single encoding technique for all types of content to achieve simplicity in design.
\\[0.1cm]
\textit{Explanation:} 
Depending on the type of content within a video, it is better to use different encoding schemes to get optimal results.
 \item Encode the entire video based on the content inside it---for example, encode an action movie with dynamic colors differently than a 2D cartoon.
\\[0.1cm]
\textit{Explanation:} 
Encoding entire video helps but it is not the most optimal solution.
 \item \textbf{[CORRECT]} Use different encoding techniques for different shots or segments within a video. For example, highly dynamic scenes should be encoded differently than lesser dynamic scenes within a single video.
\\[0.1cm]
\textit{Explanation:} 
Using a per-shot encoding technique can save storage space, transmission bandwidth, as well as make streaming convenient in terms of user bandwidth.
 \item Encoding doesn't matter at all.
\\[0.1cm]
\textit{Explanation:} 
This isn't true. Encoding helps in the overall optimization.
\end{itemize}
\vspace{0.3cm}
\noindent
\textbf{Question 4:} \textbf{Adaptive bitrate streaming} is a common technique used to provide a good streaming experience to the user. It works by detecting a user's bandwidth and the CPU capacity of the user's device in real time, and adjusts the quality of the media stream accordingly.
\\
Considering the above explanation, which diagram explains adaptive bitrate the best?
\\[0.2cm]
\begin{itemize}
 \item \includegraphics[width=\linewidth,height=3.64583in,keepaspectratio]{/api/collection/10370001/4941429335392256/page/6669179066777600/image/4546039853219840?page_type=collection_lesson&get_optimised=true}
\\[0.1cm]
\textit{Explanation:} 
Webservers don't generate multiple video qualities. It is the job of encoders.
 \item \includegraphics[width=\linewidth,height=3.64583in,keepaspectratio]{/api/collection/10370001/4941429335392256/page/6669179066777600/image/4570361581928448?page_type=collection_lesson&get_optimised=true}
\\[0.1cm]
\textit{Explanation:} 
CDNs don't generate multiple video qualities. It is the job of encoders.
 \item \includegraphics[width=\linewidth,height=3.64583in,keepaspectratio]{/api/collection/10370001/4941429335392256/page/6669179066777600/image/4906648931270656?page_type=collection_lesson&get_optimised=true}
\\[0.1cm]
\textit{Explanation:} 
Application servers run the business or logic of the application instead of serving videos to the clients.
 \item \textbf{[CORRECT]} \includegraphics[width=\linewidth,height=3.64583in,keepaspectratio]{/api/collection/10370001/4941429335392256/page/6669179066777600/image/5362086323814400?page_type=collection_lesson&get_optimised=true}
\\[0.1cm]
\textit{Explanation:} 
The encoder converts videos into multiple formats and the streaming server sends it to the client.
\end{itemize}
\vspace{0.3cm}
\noindent
\textbf{Question 5:} What is the most appropriate sequence from development to end user delivery to stream content?
\\[0.2cm]
\begin{itemize}
 \item \textbf{[CORRECT]} Encoding, bitrate conversion, subtitles addition, and deployment
 \item Deployment, bitrate conversion, subtitles, and encoding
\\[0.1cm]
\textit{Explanation:} 
It is non-sensical to first deploy a video and then encode it.
 \item Encoding, bitrate conversion, uploading to CDN, and deployment
\\[0.1cm]
\textit{Explanation:} 
Uploading to CDN is the same thing as deployment. Also, subtitles addition is missed out.
 \item Uploading to CDN, encoders, and bitrate conversion
\\[0.1cm]
\textit{Explanation:} 
First, videos must be encoded and then deployed to CDNs.
\end{itemize}
\vspace{0.3cm}
\noindent
\textbf{Question 6:} Assume that 100,000 users try to watch a football stream at 1080p in a densely populated city. If the CDN point of presence (PoP) in that area can provide a bandwidth of 300 Gbps for end users, which design best serves the clients without any buffering issues?
\\[0.2cm]
\begin{itemize}
 \item \includegraphics[width=\linewidth,height=1.97917in,keepaspectratio]{/api/collection/10370001/4941429335392256/page/6669179066777600/image/5592016216260608?page_type=collection_lesson&get_optimised=true
}
\\[0.1cm]
\textit{Explanation:} 
Not the best choice because it will require an extraordinary bandwidth for the CDN to serve out 100,000 users each with 1080p quality of streaming service.
 \item \includegraphics[width=\linewidth,height=1.97917in,keepaspectratio]{/api/collection/10370001/4941429335392256/page/6669179066777600/image/5317169078730752?page_type=collection_lesson&get_optimised=true
}
\\[0.1cm]
\textit{Explanation:} 
This strategy will put huge burden on the backbone Internet. Bottlenecks will create issues.
 \item \textbf{[CORRECT]} \includegraphics[width=\linewidth,height=1.97917in,keepaspectratio]{/api/collection/10370001/4941429335392256/page/6669179066777600/image/4957626612187136?page_type=collection_lesson&get_optimised=true
}
\\[0.1cm]
\textit{Explanation:} 
Cache servers maintained at the ISP will provide coverage to all the users reducing bandwidth requirements for the CDN as well as the origin servers.
 \item \includegraphics[width=\linewidth,height=2.60417in,keepaspectratio]{/api/collection/10370001/4941429335392256/page/6669179066777600/image/5513268930412544?page_type=collection_lesson&get_optimised=true
}
\\[0.1cm]
\textit{Explanation:} 
Streaming servers are an alternative to colocation sites. These will not be enough for 100,000 users within the same geographical location.
\end{itemize}

% End of section content